\documentclass[10pt, aspectratio=169]{beamer}
\usepackage{../beamer_theme_408}
\title{408考研零基础 --- 导学篇}
\subtitle{0-1 408考研全景解读}
\author{}
\date{}

\begin{document}

\begin{frame}{本节目标}
    \begin{enumerate}
        \item 了解 408 考试的基本概念、历史与定位
        \item 掌握四门考试科目的分值分布与难度等级
        \item 理解 408 考试的题型特点与考查方式
    \end{enumerate}
\end{frame}

\begin{frame}{什么是 408}
    \textbf{408}：计算机学科专业基础综合全国统一考试科目代码
    \vspace{0.5em}
    \begin{itemize}
        \item 2009 年起由教育部考试中心推出
        \item 涵盖四门核心课程
        \item 满分 150 分，考试时间 180 分钟
    \end{itemize}
    \vspace{0.5em}
    全国上百所高校采用，包括清华、北大、浙大、上交等
\end{frame}

\begin{frame}{考试科目与分值}
    \begin{center}
    \begin{tabular}{|c|c|c|}
        \hline
        \textbf{科目} & \textbf{分值} & \textbf{难度} \\
        \hline
        数据结构 & 45 分 (30\%) & $\bigstar$$\bigstar$$\bigstar$ \\
        \hline
        计算机组成原理 & 45 分 (30\%) & $\bigstar$$\bigstar$$\bigstar$$\bigstar$$\bigstar$ \\
        \hline
        操作系统 & 35 分 (23.3\%) & $\bigstar$$\bigstar$$\bigstar$$\bigstar$ \\
        \hline
        计算机网络 & 25 分 (16.7\%) & $\bigstar$$\bigstar$$\bigstar$ \\
        \hline
        \textbf{合计} & \textbf{150 分} & \\
        \hline
    \end{tabular}
    \end{center}
\end{frame}

\begin{frame}{考试题型}
    \textbf{单项选择题（40 道）}
    \begin{itemize}
        \item 每题 2 分，合计 80 分
        \item 数据结构约 11 道，组成原理约 11 道
        \item 操作系统约 10 道，计网约 8 道
    \end{itemize}
    \vspace{0.3em}
    \textbf{综合应用题（7 道）}
    \begin{itemize}
        \item 合计 70 分
        \item 数据结构 2 道，组成原理 2 道
        \item 操作系统 2 道，计网 1 道
    \end{itemize}
\end{frame}

\begin{frame}{备考建议}
    \textbf{学习顺序}
    \begin{enumerate}
        \item 数据结构（基础）
        \item 计算机组成原理（硬件核心）
        \item 操作系统（软硬件桥梁）
        \item 计算机网络（相对独立）
    \end{enumerate}
    \vspace{0.5em}
    \textbf{教材推荐}
    \begin{itemize}
        \item 严蔚敏《数据结构》、唐朔飞《组成原理》
        \item 汤子瀛《操作系统》、谢希仁《计网》
        \item 王道/天勤考研辅导书
    \end{itemize}
\end{frame}

\begin{frame}{时间规划}
    \begin{tabular}{|c|c|c|}
        \hline
        \textbf{阶段} & \textbf{时长} & \textbf{任务} \\ \hline
        零基础阶段 & 2-3 个月 & 本系列课程 + 知识框架 \\ \hline
        强化阶段 & 3-4 个月 & 教材精读 + 王道刷题 \\ \hline
        冲刺阶段 & 1-2 个月 & 模拟题 + 查漏补缺 \\ \hline
    \end{tabular}
\end{frame}

\begin{frame}{易错点：408常见误解}
    \begin{alertblock}{误解一：408只考数据结构与算法}
        事实：408包含数据结构(45分)、组成原理(45分)、操作系统(35分)、计网(25分)，四科缺一不可。
    \end{alertblock}
    \vspace{0.3em}
    \begin{alertblock}{误解二：组成原理可以突击学习}
        事实：组成原理是408中\textbf{难度最大}的科目，抽象程度高，需要长期积累理解，建议安排最多学习时间。
    \end{alertblock}
    \vspace{0.3em}
    \begin{alertblock}{误解三：计网分少可以放弃}
        事实：计网25分虽然占比最小，但知识点相对独立、规律性强，是\textbf{性价比最高}的科目，认真复习容易拿分。
    \end{alertblock}
\end{frame}

\begin{frame}{课堂练习}
    \textbf{检验对本节内容的掌握：}
    \vspace{0.5em}
    \begin{enumerate}
        \item 408考试的四门科目是什么？各自的分值是多少？
        \item 408备考可以分为哪三个阶段？各阶段的主要任务是什么？
    \end{enumerate}
    \vspace{0.5em}
    \begin{block}{参考答案}
        \begin{enumerate}
            \item 数据结构(45分)、计算机组成原理(45分)、操作系统(35分)、计算机网络(25分)
            \item 零基础阶段(2-3月)：建立知识框架；强化阶段(3-4月)：教材精读+刷题；冲刺阶段(1-2月)：模考+查漏补缺
        \end{enumerate}
    \end{block}
\end{frame}

\begin{frame}{本节总结}
    \begin{enumerate}
        \item 408 是计算机考研统考科目代码
        \item 数据结构 45 分 + 组成原理 45 分 + 操作系统 35 分 + 计网 25 分
        \item 题型：40 道单选 + 7 道综合题
        \item 推荐学习顺序：数据结构 $\to$ 组成原理 $\to$ 操作系统 $\to$ 计网
    \end{enumerate}
\end{frame}

\begin{frame}{下节预告}
    \begin{center}
    \vspace{1em}
    {\Large \textbf{0-2 计算机系统初体验}}\\
    \vspace{0.3em}
    \textcolor{408blue}{从按下电源到操作系统加载的完整过程}
    \vspace{1em}
    \end{center}
    \begin{block}{下节内容}
        \begin{itemize}
            \item 计算机系统的层次结构
            \item 五大核心部件（冯·诺依曼体系）
            \item 从按下电源到桌面的启动六步骤
            \item 软件与硬件的关系
        \end{itemize}
    \end{block}
    \vspace{0.5em}
    \begin{center}
        \textcolor{408orange}{敬请期待！}
    \end{center}
\end{frame}
\end{document}
